\section{Discussion and conclusion}
\label{sec:discussion}

\subsection{What the cases show}
\label{sec:discussion-findings}

AERead evaluates an agent's economic decisions by the value it obtains for
a principal, given the information, resources and authority of its seat.
Each score is measured against a reference whose access to information is
stated, whether a feasible policy, the best policy under the seat's
information or a full-information bound, and it keeps its own units and
names the counterparts it was measured against, so that every shortfall
has a defined interpretation. The worked examples apply this to
information acquisition, bargaining and commitment in housing, procurement
and contracting, and a reimplementation of part of TERMS-Bench runs on the
same runtime.

The cases teach three lessons that bear on any evaluation of this kind.
The counterpart is part of the measurement: the same two agents differ by
far more when they negotiate with each other than when each faces a fixed
counterpart. The brief is part of the case: stating a rule it had left
implicit moved scores as much as changing the agent did. And the defects
that matter are often properties of the population of worlds, the brief or
the analysis rather than of any single world, which is why construction
and quality control belong in the method; the checks of Section 6 changed
the interpretation of our own results. Loss categories locate a shortfall
but do not explain it, and explaining it requires controlled follow-up
comparisons. Replay establishes consistency with recorded play; construct
validity requires evidence about the worlds, the information and the
controls.

These properties matter most where agents are deployed. Organizations that
pilot agents often stop because the business value is unclear and the
risks are hard to control \citep{gartner2025agentic}; a value score against
a stated reference and a typed account of losses address both. Small firms
face a different problem: short of staff \citep{nfib2026jobs}, they often
leave decisions like those our cases model to one person, without
specialized expertise and with no review before a commitment. Delegating
such a decision to an agent is attractive there, and trusting it requires
evidence that the agent keeps value against the counterparts it will meet.

The present evidence has clear limits. The cases take their decision
structure from informal rentals, wholesale sourcing and AI-infrastructure
contracting, but their parameters are set by the authors; calibration to
field data and transfer to field settings remain to be established.
Coverage is limited to two transaction classes and a small set of agent
configurations, with no human policy baseline. Housing v15 and Procurement's
declared pack pass four of the five qualification gates and no case passes
all five: several checks were added after the runs they qualify, and hosted-model and CLI
behavior limits attribution from declared settings.
Versioned reanalyses now regenerate the corrected Housing measures and
Procurement decision losses from sealed records, which strengthens their
provenance, but full verification against original execution still
requires retained logs (Appendix~G), and the reanalyses neither qualify the
campaigns retroactively nor resolve Housing's historical seat attribution
(Appendix~H). The results therefore compare named agents on stated worlds;
they do not establish general model rankings or deployment readiness.

\subsection{Open research directions and contributions}
\label{sec:open-directions}

AERead is shared infrastructure for studying how agents make economic
decisions. Its code is open source under the Apache-2.0 license
(github.com/aeread-org/AERead; contact@aeread.org), and contributors can
pursue their own questions with the existing cases, counterparts and
evidence. The results above leave four questions open, each suited to a
different kind of contribution.
\begin{itemize}
\item \textbf{Counterpart populations.} The Data-center gap between two
agents is 42\% of the value available in a deal head to head and 7\%
against a fixed script.
How far does an agent's measured skill depend on whom it faces, can
rankings reverse across counterpart populations (none did in our runs),
and which populations separate distinct strategic capabilities? The
scripted styles, head-to-head pairings and counterpart seats are in place.
\item \textbf{Validity of populations of worlds.} Of 25 defects selected
from our record, 14 were invisible in any single world. Which checks find
such defects before a campaign runs, and what should qualify a case? The
quality-control gates, the incident record and the failed campaigns are
available for study.
\item \textbf{Information and judgment.} Stating a rule the brief had left
implicit moved scores as much as changing the agent. Separating what an
agent was told from how it decided, and computing best-informed references
for richer cases, are open methodological problems.
\item \textbf{New economic decisions.} A case for another activity
supplies its rules, worlds and reference; execution, replay, counterpart
seats and the quality-control standard come with the runtime.
\end{itemize}

What is ready and what is not: the runtime, its sealed evidence and
replay, and the case layout are in use, and Housing, Procurement and Data
center each have a world generator, references and scripted counterparts.
No case is yet fully qualified: Housing v15 and Procurement's declared pack
pass four of the five gates.
Section~\ref{sec:qc} gives every gate verdict and
Appendix~F each campaign's open checks, so a contributor can see what a
case still needs before its results can rank agents.

The next steps follow from the limits in Section 7.1. Future campaigns should bind
actual seat assignments and headline scorers into frozen evaluation
identities and run the required controls before any comparison, and the
TERMS-Bench reimplementation needs the protocol fixes described in
Section 5.3 before a declared reproduction. Practitioner review would
calibrate the cases' parameters, human baselines would show where agents
stand against people, and externally contributed cases would test whether
the structure serves authors other than its maintainers. AERead's present
contribution is that structure, a common way to construct evaluations of
economic decisions, together with an evidence-based account of the checks
needed to interpret what they measure.
